\documentclass[../DoAn.tex]{subfiles}
\begin{document}

Sự phát triển mạnh mẽ của Công nghệ Thông tin (CNTT) trong hai thập kỷ qua đã
và đang làm thay đổi sâu rộng phương thức vận hành của hầu hết các lĩnh vực
trong đời sống xã hội, từ kinh tế, y tế, công nghiệp cho đến giáo dục và quản
lý công. Việc ứng dụng các giải pháp tin học không chỉ dừng lại ở mức tự động
hoá các tác vụ thủ công lặp lại, mà còn mở ra khả năng phân tích, dự báo và hỗ
trợ ra quyết định trong các hệ thống phức tạp vốn đòi hỏi sự cân bằng giữa
nhiều ràng buộc đối lập nhau. Trong bối cảnh đào tạo đại học, xu hướng số hoá
đang dần chuyển hoá không chỉ hoạt động dạy và học mà còn cả các quy trình hành
chính và dịch vụ sinh viên, trong đó quản lý ký túc xá là một trong những mảng
đang đòi hỏi sự cải tiến rõ rệt nhất.

Ký túc xá đại học không đơn thuần là nơi cung cấp chỗ ở; đây là một phần quan
trọng trong hệ sinh thái hỗ trợ học tập, ảnh hưởng trực tiếp đến chất lượng
cuộc sống, năng lực tập trung học tập và sự hài lòng của sinh viên trong suốt
quá trình học tập tại trường. Một hệ thống quản lý ký túc xá hoạt động tốt cần
đảm bảo tính công bằng trong phân bổ phòng, tính minh bạch trong chính sách ưu
tiên, tính kịp thời trong thông báo kết quả, và tính linh hoạt khi chính sách
nhà trường thay đổi. Những yêu cầu này, khi được giải quyết bằng phương pháp
thủ công truyền thống, thường dẫn đến sai sót, chậm trễ và thiếu nhất quán,
đặc biệt trong bối cảnh quy mô sinh viên ngày càng tăng và cấu trúc chính sách
ngày càng phức tạp.

Xuất phát từ thực tiễn trên, đề tài này tập trung nghiên cứu và phát triển một
hệ thống phần mềm tích hợp phục vụ quản lý và phân bổ ký túc xá, kết hợp giữa
một Allocation Engine có khả năng cấu hình chính sách linh hoạt với giao diện
người dùng hiện đại được tối ưu hoá cho thiết bị di động. Hệ thống hướng tới
việc giải quyết đồng thời ba nhóm vấn đề cốt lõi: tự động hoá và tối ưu hoá
quy trình phân bổ; đảm bảo tính minh bạch và khả năng kiểm chứng của mọi quyết
định; và nâng cao trải nghiệm người dùng thông qua giao tiếp thời gian thực và
trực quan hoá không gian. Chương 1 được tổ chức thành bốn phần: đặt vấn đề
(1.1), mục tiêu và phạm vi nghiên cứu (1.2), định hướng giải pháp (1.3) và bố
cục đồ án (1.4).

\section{Đặt vấn đề}
\label{section:1.1}

Bài toán phân bổ phòng ở ký túc xá mang đặc trưng của một bài toán tối ưu hoá
có ràng buộc đa chiều, trong đó mỗi quyết định gán phòng phải đồng thời thoả
mãn nhiều điều kiện: chính sách ưu tiên của nhà trường, các ràng buộc về sức
chứa và an toàn, tính công bằng giữa các nhóm sinh viên khác nhau, và yêu cầu
minh bạch để người dùng có thể kiểm tra và khiếu nại khi cần. Độ phức tạp của
bài toán tăng lên đáng kể khi số lượng ứng viên lớn, khi chính sách có nhiều
tầng ưu tiên chồng nhau, hoặc khi cần xử lý các trường hợp đặc biệt như sinh
viên có hoàn cảnh khó khăn, sinh viên trao đổi quốc tế, hay sinh viên có nhu
cầu đặc thù về sức khoẻ.

Thực trạng hiện nay tại Đại học Bách Khoa Hà Nội cũng như tại nhiều trường đại
học quy mô lớn trong nước cho thấy quy trình phân bổ phòng ký túc xá vẫn chủ
yếu dựa vào bảng tính điện tử hoặc xử lý thủ công qua nhiều bước nối tiếp: thu
thập hồ sơ đăng ký, tổng hợp thông tin ưu tiên, tính điểm theo quy định, xếp
hạng ứng viên, gán phòng theo thứ tự và kiểm tra lại để đảm bảo không vi phạm
ràng buộc. Quy trình này không chỉ tiêu tốn nhiều công sức của cán bộ mà còn
tiềm ẩn nhiều rủi ro: tính toán nhầm do thao tác tay, sơ sót dữ liệu khi hợp
nhất từ nhiều nguồn, và hiểu sai chính sách khi quy định được diễn giải khác
nhau bởi những người thực hiện khác nhau. Hậu quả là kết quả phân bổ đôi khi
không nhất quán, khó truy nguyên, và mất nhiều thời gian để xử lý khi có khiếu
nại.

Thiếu minh bạch là một trong những hạn chế nghiêm trọng nhất của hệ thống hiện
tại. Sinh viên không có cách nào kiểm chứng cách điểm ưu tiên của mình được
tính toán, vì sao một sinh viên khác lại được ưu tiên hơn, hay liệu kết quả phân
bổ có tuân thủ đúng chính sách đã được công bố hay không. Sự thiếu minh bạch
này dẫn đến mất lòng tin, gia tăng số lượng khiếu nại và buộc cán bộ phải dành
đáng kể thời gian để giải thích và xử lý tranh chấp thay vì tập trung vào công
tác chuyên môn.

Bên cạnh đó, hạ tầng dữ liệu còn nhiều bất cập: thông tin sinh viên và thông
tin phòng thường lưu trữ ở các hệ thống rời rạc với định dạng không chuẩn hoá,
khiến việc đồng bộ và đối chiếu tốn công sức. Khi chính sách thay đổi, cán bộ
phải thực hiện lại toàn bộ quy trình thủ công trên tập dữ liệu mới, không có
cách nào mô phỏng trước tác động của chính sách mới để so sánh các phương án.
Việc thiếu nhật ký thao tác (audit log) có cấu trúc cũng làm cho quá trình truy
nguyên nguyên nhân sai sót trở nên cực kỳ khó khăn.

Cuối cùng, kỳ vọng của sinh viên về trải nghiệm số ngày càng cao. Họ muốn đăng
ký và theo dõi trạng thái hồ sơ từ điện thoại di động, xem trước ảnh 360° hoặc
sơ đồ không gian phòng trước khi đưa ra quyết định, nộp minh chứng điện tử và
nhận thông báo tức thì khi có kết quả. Những nhu cầu này hoàn toàn nằm ngoài
khả năng đáp ứng của quy trình thủ công hiện tại.

Tóm lại, bài toán cần giải quyết bao gồm sáu nhóm vấn đề có liên quan chặt chẽ
với nhau: (1)~quy trình phân bổ thủ công, tiêu tốn thời gian và tiềm ẩn sai sót;
(2)~thiếu minh bạch và khả năng kiểm chứng trong toàn bộ quy trình; (3)~dữ liệu
phân tán, thiếu chuẩn hoá và khó bảo trì theo thời gian; (4)~không có khả năng
mô phỏng và đánh giá tác động trước khi áp dụng chính sách; (5)~giao diện chưa
tối ưu cho thiết bị di động và nhu cầu trải nghiệm số của sinh viên; và
(6)~thiếu cơ chế phản hồi kịp thời, minh bạch giữa hệ thống và người dùng cuối.

\section{Mục tiêu và phạm vi đề tài}
\label{section:1.2}

\subsection{Mục tiêu nghiên cứu}

Mục tiêu tổng quát của đề tài là nghiên cứu, thiết kế và triển khai một hệ
thống phần mềm quản lý ký túc xá tích hợp, trong đó cốt lõi là một Allocation
Engine có khả năng thực thi chính sách phân bổ linh hoạt, kết hợp với các giao
diện người dùng hiện đại và cơ chế đảm bảo minh bạch. Hệ thống hướng đến việc
thay thế hoàn toàn quy trình thủ công hiện tại, đồng thời cung cấp nền tảng kỹ
thuật có thể mở rộng và bảo trì lâu dài.

Về phía Allocation Engine, hệ thống cần thiết kế một thành phần lõi cho phép
mô tả chính sách phân bổ dưới dạng dữ liệu có cấu trúc, hỗ trợ hai chế độ
hoạt động là mô phỏng (preview) và thực thi (commit), đồng thời tự động sinh
nhật ký giải thích cơ sở quyết định cho từng sinh viên để đảm bảo tính
auditability. Về phía backend và API, cần xây dựng một tầng dịch vụ an toàn và
có khả năng mở rộng, cung cấp API theo chuẩn RESTful được bảo vệ bằng xác thực
và phân quyền, lưu trữ lịch sử phân bổ theo phiên và ghi đầy đủ audit log cho
các thao tác quan trọng.

Về phía giao diện quản trị, cần phát triển giao diện cho cán bộ quản lý nhằm
cấu hình chính sách, khởi chạy và giám sát phiên phân bổ, điều chỉnh kết quả
thủ công khi cần và xuất báo cáo thống kê. Về phía giao diện sinh viên, cần
thiết kế trải nghiệm mobile-first cho phép sinh viên đăng ký, nộp minh chứng
điện tử, tra cứu điểm ưu tiên với lời giải thích có thể kiểm chứng, và nhận
thông báo kết quả kịp thời. Ngoài ra, hệ thống cần tích hợp cơ chế thông báo
thời gian thực thông qua WebSocket, cùng với bộ kiểm thử bao gồm kiểm thử đơn
vị, tích hợp và tải nhằm đảm bảo độ chính xác và hiệu năng trong điều kiện vận
hành thực tế.

\subsection{Phạm vi đề tài}

Phạm vi triển khai bao gồm toàn bộ các thành phần tạo nên một hệ thống hoàn
chỉnh có thể vận hành độc lập: tầng backend và API được xây dựng trên nền
Node.js/Express với MongoDB Atlas làm cơ sở dữ liệu và Redis làm tầng hỗ trợ
realtime; Allocation Engine phía máy chủ với đầy đủ chế độ mô phỏng và thực
thi; giao diện quản trị và sinh viên render phía máy chủ sử dụng EJS kết hợp
Bootstrap; ứng dụng di động đa nền tảng với Expo/React Native; hệ thống thông
báo thời gian thực với Socket.IO; cơ chế xác thực và phân quyền dựa trên JWT
và RBAC; cùng bộ kiểm thử và tài liệu hướng dẫn triển khai đầy đủ.

Về phạm vi thiết bị và khả năng tiếp cận, hệ thống được xây dựng theo nguyên
tắc responsive nhằm đảm bảo trải nghiệm nhất quán trên bốn nhóm thiết bị: máy
tính để bàn, laptop, máy tính bảng và điện thoại thông minh. Giao diện web tự
động điều chỉnh bố cục theo chiều rộng màn hình thông qua CSS Grid, trong khi
ứng dụng di động React Native cung cấp trải nghiệm native trên cả iOS và Android.
Điều này phản ánh xu hướng tất yếu trong thiết kế hệ thống giáo dục hiện đại:
sinh viên truy cập dịch vụ từ nhiều loại thiết bị khác nhau và kỳ vọng vào một
trải nghiệm thống nhất, không phụ thuộc vào nền tảng.

Bên cạnh đó, hệ thống được thiết kế có tính đến yêu cầu hỗ trợ đa ngôn ngữ,
phục vụ không chỉ sinh viên trong nước mà còn nhóm sinh viên quốc tế, sinh viên
trao đổi và du học sinh. Trong bối cảnh Đại học Bách Khoa Hà Nội và các trường
đại học kỹ thuật tại Việt Nam đang mở rộng hợp tác quốc tế, đón nhận ngày càng
nhiều sinh viên từ các quốc gia trong khu vực và trên thế giới, nhu cầu sử dụng
hệ thống ký túc xá không còn giới hạn trong nhóm sinh viên nói tiếng Việt. Một
hệ thống quản lý ký túc xá hiện đại cần đảm bảo khả năng tiếp cận rộng hơn,
trong đó rào cản ngôn ngữ phải được xem xét từ giai đoạn thiết kế ban đầu, thay
vì bổ sung tùy hứng về sau.

Một số nội dung được xác định là nằm ngoài phạm vi của đề tài này: tích hợp
sâu với các hệ thống nội bộ của trường như hệ thống quản lý điểm hay tuyển sinh
(chỉ tích hợp qua API mẫu); phát triển các mô-đun phân tích dữ liệu nâng cao
hay mô hình dự báo bằng học máy; và triển khai hạ tầng vận hành theo tiêu chí
sản xuất toàn diện với cân bằng tải đa vùng, sao lưu địa lý và khôi phục thảm
họa. Những nội dung này có thể được phát triển trong các giai đoạn tiếp theo khi
hệ thống được đưa vào vận hành thực tế.

\section{Định hướng giải pháp}
\label{section:1.3}

Giải pháp được thiết kế xung quanh bốn nguyên tắc cốt lõi: phân tách trách
nhiệm rõ ràng giữa các tầng; khả năng cấu hình chính sách mà không cần thay
đổi mã nguồn; ưu tiên trải nghiệm người dùng trên thiết bị di động; và đảm bảo
khả năng mô phỏng, kiểm thử trước khi áp dụng chính sách vào thực tế.

Về kiến trúc hệ thống, hệ thống được tổ chức theo mô hình ba tầng với phân tách
rõ ràng về trách nhiệm: tầng trình bày xử lý mọi tương tác người dùng; tầng
dịch vụ chứa logic nghiệp vụ bao gồm Allocation Engine; và tầng dữ liệu quản
lý lưu trữ và truy vấn. Allocation Engine tiếp nhận đầu vào là danh sách ứng
viên, danh sách phòng và cấu hình chính sách, sau đó tính toán điểm ưu tiên
theo các tiêu chí đã thiết lập, xếp hạng và tiến hành gán phòng theo thứ tự ưu
tiên trong khi tôn trọng các ràng buộc cứng về sức chứa và an toàn. Sự phân
tách giữa chế độ mô phỏng và chế độ thực thi cho phép cán bộ đánh giá toàn
diện tác động của mỗi kịch bản chính sách trước khi ra quyết định chính thức.

Về lựa chọn công nghệ, tập hợp công nghệ được lựa chọn nhằm cân bằng giữa tốc
độ phát triển, hiệu năng vận hành và tính ổn định của hệ sinh thái thư viện.
Node.js với Express được chọn cho backend bởi khả năng xử lý bất đồng bộ tự
nhiên và hệ sinh thái npm phong phú. MongoDB Atlas được sử dụng như cơ sở dữ
liệu đám mây chính nhờ mô hình document linh hoạt, phù hợp với cấu trúc chính
sách đa dạng và có thể thay đổi. Redis đảm nhiệm tầng hỗ trợ cho Socket.IO
cluster adapter, cho phép mở rộng ngang trong tương lai. Phía giao diện, cả
admin portal và student portal đều sử dụng EJS để render phía máy chủ kết hợp
Bootstrap; ứng dụng di động được phát triển với Expo và React Native để tái sử
dụng mã nguồn trên cả iOS và Android.

Về bảo mật và phân quyền, hệ thống áp dụng chiến lược bảo mật đa tầng: web
portal sử dụng session-based authentication với cookie HttpOnly và SameSite;
mobile API sử dụng JWT access token (TTL 15 phút) kết hợp refresh token (TTL 30
ngày) với cơ chế phát hiện anomaly. QR thẻ cư trú được ký số bằng HMAC-SHA256.
Toàn bộ API được bảo vệ bằng kiểm soát truy cập theo vai trò (RBAC). Xác thực
hai yếu tố TOTP được áp dụng cho tài khoản quản trị. Môi trường vận hành được
kiểm tra toàn bộ biến cấu hình trước khi khởi động.

Về kiểm thử và đảm bảo chất lượng, chiến lược kiểm thử được thiết kế theo
nhiều tầng: kiểm thử đơn vị cho Allocation Engine và các service lõi; kiểm thử
tích hợp cho các API chính; kiểm thử giao diện người dùng; và kiểm thử tải
đánh giá khả năng đáp ứng trong điều kiện vận hành thực tế.

\section{Bố cục đồ án}
\label{section:1.4}

Phần còn lại của báo cáo đồ án tốt nghiệp được tổ chức thành năm chương, mỗi
chương có nhiệm vụ rõ ràng và liên kết logic với nhau, tạo thành một chuỗi lập
luận nhất quán từ cơ sở lý thuyết đến triển khai và đánh giá.

Chương 2 --- Cơ sở lý thuyết và công trình liên quan khảo sát và phân tích các
nền tảng lý thuyết liên quan đến bài toán phân bổ tài nguyên, các thuật toán
xếp hạng và ghép cặp được sử dụng phổ biến trong nghiên cứu, cùng với tổng
quan các hệ thống quản lý ký túc xá đang tồn tại để xác định khoảng trống mà
đề tài hướng tới lấp đầy.

Chương 3 --- Công nghệ và kiến trúc hệ thống trình bày đầy đủ các yêu cầu chức
năng và phi chức năng được rút ra từ khảo sát thực tiễn, phân tích các trường
hợp sử dụng theo khuôn mẫu học thuật, thiết kế kiến trúc hệ thống tổng thể và
mô hình dữ liệu, kết thúc bằng thiết kế chi tiết Allocation Engine bao gồm
thuật toán và cấu trúc dữ liệu.

Chương 4 --- Thiết kế, triển khai và đánh giá hệ thống mô tả các quyết định
kiến trúc phần mềm, cấu trúc package của dự án, chi tiết triển khai các thành
phần cốt lõi bao gồm Allocation Engine, hệ thống realtime, bảo mật đa tầng và
ứng dụng di động, cùng hướng dẫn triển khai môi trường hoàn chỉnh.

Chương 5 --- Kiểm thử và đánh giá trình bày chiến lược kiểm thử toàn diện, mô
tả các kịch bản kiểm thử chức năng, tích hợp và hiệu năng, phân tích kết quả
thu được và đánh giá mức độ đáp ứng yêu cầu đề ra ở Chương~3.

Chương 6 --- Kết luận và hướng phát triển tương lai tổng kết các kết quả đạt
được, nêu rõ đóng góp chính của đồ án, nhận diện các hạn chế còn tồn tại và
đề xuất hướng phát triển tiếp theo trong các giai đoạn sau khi hệ thống được
đưa vào vận hành thực tế.

Ngoài sáu chương chính, báo cáo bao gồm phần Tài liệu tham khảo liệt kê các
bài báo khoa học, sách chuyên khảo, thông số kỹ thuật chính thức và tài liệu
từ các dự án mã nguồn mở có liên quan. Phần Phụ lục cung cấp các thông tin bổ
sung kỹ thuật bao gồm tài liệu API đầy đủ, sơ đồ cơ sở dữ liệu, hướng dẫn cài
đặt và triển khai, script kiểm thử và các kịch bản thử nghiệm, nhằm hỗ trợ các
nhà phát triển muốn hiểu sâu, tái triển khai hoặc mở rộng hệ thống trong tương
lai.

Như vậy, chương 1 đã phác thảo bức tranh tổng quan về bài toán quản lý và phân
bổ ký túc xá, từ thực trạng tồn tại đến mục tiêu cần đạt được và định hướng
giải pháp. Cơ sở nhận thức được xây dựng trong chương này làm tiền đề để chương
2 đi sâu vào khảo sát thực tế và nghiên cứu các hệ thống tương tự, từ đó củng
cố cơ sở lý luận cho thiết kế và triển khai hệ thống ở các chương tiếp theo.

\end{document}
